\section{Discussion}

\subsection{Implications for intelligent-interface authoring}

The explicit path gives developers a pre-interaction view of responsibility,
permitted behavior, expected outcomes, and failures such as unknown targets,
ambiguous owners, missing bindings, or contradictory observations. This is a
diagnostic claim, not an explanation of model reasoning, and a static
single-chatbot interface may not justify the added structure. The workbench
currently revises placement only. Responsibility, capability, and handoff
relations are a roadmap: they should become typed preview--validate--commit
transactions rather than unrestricted graph edits.

\subsection{Control and failure recovery}

Silent fallback after failed grounding can produce a fluent but misleading
response. The visitor path instead exposes no-target and ambiguity status and
blocks the answer transition; route and assertion details remain
developer-facing. Clarification, transparent-handoff, and abstention dialogues
are future interaction work.

\subsection{Composing interaction and authorization policy}

``Authorized'' is local to the interaction model: the scenario declares this
provider and capability for this target. It neither authenticates a person nor
grants access to a protected resource. A deployment can resolve the exact
target--provider--action tuple with \systemname{} and then require an external
relationship- or capability-based policy to approve it; lifecycle governance
can separately decide which implementation may be active. The preflight
boundary permits this composition, but the prototype integrates neither
mechanism.

\subsection{Practical and societal considerations}

Spatial-agent traces can expose rooms, locations, speech, and interaction
histories. The backend receives the utterance, target identifiers, ray
metadata, and modality; the model provider receives bounded history and trusted
entity/group/zone context but not raw ray coordinates. Persistent traces omit
utterance text, yet the prototype lacks visitor-facing consent, retention,
export, and deletion controls. Deployments should minimize data and separate
technical identifiers from personal identities. Agent roles can also encode
organizational bias: inspectability does not replace review of representation,
authority, refusal, safeguards, or knowledge provenance.

\subsection{Artifact-release boundary}

The review artifact supports local replication but grants no blanket
redistribution license. Public release requires an owner-approved license, a
third-party inventory, and removal or replacement of restricted Unity/media
assets. The intended path is a clean tagged snapshot plus a newly licensed
minimal scene that runs the same contract tests.

\subsection{When a model helps---and when it does not}

The approach trades authoring structure for evidence and helps only while model
and runtime remain synchronized. Both adapters detected the common faults;
model-grounded patches changed fewer artifacts but more references and ran
slower locally. Whether this trade helps authors requires human evaluation.
